Claims about where a language model stores in-context state, and about which component carries an intervention's effect, are usually measured under one answer format. Later tokens read a context token only through its attention key (whether they attend to it) and its value (what they copy). We clamp these channels separately and exactly at the token where a story writes an object's location, under five readouts that differ only in whether and where the candidate answers are listed and in the matching answer instruction, in \nModels{} models (1.5B--72B, four families), preregistering predictions for the confirmatory runs and reporting those that failed. When the candidates are listed after the state token, later tokens look the state up through its key: keys carry \shareOlmoSeven{}--\shareMistralSeven{} of the multiple-choice effect in models of 7B and above (\shareQwenOnefive{} at 1.5B), and option-word rows alone recover \locOptChoicewords{} of the key effect at 1.5B. From 7B upward, a neutral sentence re-mentioning the candidates also opens the key channel, more weakly; this was not predicted. With no candidate named after the state token, the state is copied from the value; with options before the story, keys carry at most \beforeMax{} nats of identity. At Qwen2.5-72B, a learned patch from prior work shifts behaviour by nearly the same fraction in every format ($\varphi$ \phiNoneQwen{}--\phiBeforeQwen{}), yet adding its keys alone to a PCA patch recovers \psiOptQwen{} of their difference when options follow the question and \psiNoneQwen{} when nothing re-mentions the candidates. A preregistered law predicting free-form transfer from value-channel completeness fails at Mistral-24B and holds only at the edge of its margin at Qwen2.5-72B. The component-level explanation of this intervention, and potentially of others, is readout-relative.
